\chapter{Capabilities}\label{ch:caps}

A \term{capability} is an unforgeable description of a segment of a far space and of what its holder
may do with it. Capabilities live only in capability registers, in the save table, in inboxes, and in
the hands of devices; they never appear in memory.

\section{Format}

\begin{figure}[htbp]
\centering
\begin{tikzpicture}[f/.style={draw=accent,fill=accentlight,minimum height=9mm,align=center,font=\small}]
  \node[f,minimum width=10mm] (v) at (0,0) {valid\\\scriptsize 1};
  \node[f,minimum width=11mm,right=0pt of v] (s) {space\\\scriptsize 3};
  \node[f,minimum width=27mm,right=0pt of s] (p) {perms\\\scriptsize\perm{V} \perm{D} \perm{A} \perm{X} \perm{W} \perm{R}};
  \node[f,minimum width=25mm,right=0pt of p] (b) {base\\\scriptsize 64};
  \node[f,minimum width=25mm,right=0pt of b] (l) {len\\\scriptsize 64};
  \node[f,minimum width=10mm,right=0pt of l] (d) {rv\\\scriptsize 1};
  \node[f,minimum width=12mm,right=0pt of d] (g) {gen\\\scriptsize 16};
\end{tikzpicture}
\caption{Logical contents of a capability. The in-register representation is implementation-defined.}
\label{fig:capfmt}
\end{figure}

\begin{description}[leftmargin=1em,style=nextline]
\item[valid] Whether the register or slot holds a capability. An empty register has valid = 0 and all
  other fields zero (\texttt{NULLCAP}).
\item[space] 2 (T2), 3 (T3), 4 (device), 5 (core control) or 6 (endpoint).
\item[base, len] The segment is $[\textit{base}, \textit{base} + \textit{len})$. \textit{base} is an
  absolute address in the space and always a multiple of 8 (Section~\ref{sec:capalign}); for an endpoint
  it is a node number. \textit{len} is in bytes and need not be a multiple of 8. A segment of length zero
  is valid: it permits no access but can be split, joined, given and taken like any other.
\item[perms] A set of the six permissions of Table~\ref{tab:perms}.
\item[rv] Whether the capability is \term{revocable} (Section~\ref{sec:domains}).
\item[gen] For a revocable capability, the generation of the \term{granule} that holds its base, as it
  was when the capability was made. A revocable capability whose generation is no longer its
  granule's is dead. For a revoker (\perm{V}), the permissions \insn{crevoke} and \insn{cundom} give;
  for any other capability, 0.
\end{description}

\begin{table}[htbp]
\centering\small
\caption{Capability permissions}\label{tab:perms}
\begin{tabularx}{\linewidth}{@{}c c X@{}}
\toprule
\textbf{Bit} & \textbf{Perm} & \textbf{Allows}\\\midrule
0 & \perm{R} & \insn{tget}: far $\to$ D-LS; for an endpoint, \insn{crecv}\\
1 & \perm{W} & \insn{tput}: D-LS $\to$ far; for an endpoint, \insn{csend}. Never combined with \perm{D} in memory spaces\\
2 & \perm{X} & \insn{tgetx}: far $\to$ I-LS; \insn{xmap}, \insn{xcall}\\
3 & \perm{A} & far atomics (\insn{tamo}) only\\
4 & \perm{D} & duplication: the capability is \term{copyable} rather than linear\\
5 & \perm{V} & revoker: \insn{crevoke} and \insn{cundom} over its range; grants no access\\
\bottomrule
\end{tabularx}
\end{table}

\insn{cbase}, \insn{clen} and \insn{cperms} read the fields into integer registers; \insn{cperms}
returns \texttt{space << 8 | perms}. \insn{cvalid} tests whether a capability is live. The revocable
flag and the generation cannot be read. No instruction converts an integer into a capability that grants access to
memory: capabilities can be inspected but never forged.

\section{Live capabilities}\label{sec:live}

A capability can be used only while it is \term{live}:

\begin{pseudo}
func Live(c)
    if not c.valid or not c.rv then return c.valid; end   // revokers are never revocable
    g = Granule(c.space, c.base);                         // none past the end of the space
    return g != none and GEN[g] == c.gen;
end
\end{pseudo}

Every instruction that uses a capability checks \texttt{Live} at the moment it executes and traps with
\cause{capinv} if the check fails. A capability with a stale generation is \term{dead}: it still
occupies its register or slot, but it behaves exactly like an empty one. It cannot be used,
\insn{cvalid} reports 0, and it may be overwritten like an empty register. Revocation therefore
needs no search for copies: advancing a granule's generation kills every revocable capability based
in it that was made under an older one, wherever it is.

\section{Linear and copyable capabilities}\label{sec:linear}

A capability without \perm{D} is \term{linear}. Linear capabilities obey a strict conservation law:
\begin{itemize}
\item every instruction that moves a linear capability (\insn{cmov}, \insn{crestrict}, \insn{cshare},
  \insn{catomic}, \insn{csave}, \insn{crestore}, \insn{cgive}, \insn{csend}) empties its source;
\item splitting a linear capability produces two linear capabilities for \emph{disjoint} ranges, and
  joining consumes one of its operands;
\item no instruction copies one.
\end{itemize}
A capability with \perm{D} is \term{copyable}: moving it leaves the source in place.

\begin{rule*}{Writable capabilities are linear}
No capability for T2 or T3 has both \perm{W} and \perm{D}. Because \perm{W} can never be added to a
capability (\insn{crestrict} only removes permissions, \insn{cjoin} requires equal permissions,
\insn{cshare} and \insn{catomic} drop \perm{W}), it follows that for every byte of T2 or T3, at most one
\emph{live} writable capability covering it exists anywhere in the machine.
\end{rule*}

Revocation is the one operation that creates a writable capability for memory that already had one,
and it first kills the old one: every revocable capability based in the units it touches, which are
all the capabilities that can reach the revoked range (Section~\ref{sec:domains}), so the rule holds
across revocation.

Device and core-control capabilities are exempt: the device root \reg{c3}, the console capability
\reg{c2} and the core-control capability \reg{c4} are copyable and writable, so several cores may hold
the console. Ordering output from several cores to one device is the program's responsibility
(Section~\ref{sec:determinism}). Endpoints are copyable too.

\section{No silent overwrites}\label{sec:nooverwrite}

\begin{rule*}{Live linear capabilities are never overwritten}
An instruction that would place a capability in a register or save-table slot that already holds a
live linear (non-\perm{D}) capability traps with cause \cause{capop} instead.
\end{rule*}

This applies to the destination of \insn{cmov}, \insn{crestrict}, \insn{cshare} and \insn{catomic}
(when $\reg{cd} \neq \reg{cs}$), \insn{csplit}, \insn{ctake}, \insn{crecv}, \insn{crestore},
\insn{cdom}, \insn{crevoke}, \insn{cundom}, \insn{cmkep}, and to the slot of \insn{csave}. A
destination is \term{replaceable} if it is not live or holds a copyable capability:
\begin{pseudo}
func Replaceable(c) = not Live(c) or c.D
func Linear(c)      = not c.D
\end{pseudo}
\insn{cclear} is the one explicit way to discard a capability, linear or not. The rule turns a common
bug, losing authority by overwriting a register, into a precise trap, and it makes every loss of a
linear capability visible in the program text.

\section{Deriving capabilities}

\begin{table}[htbp]
\centering\small
\caption{Capability operations}\label{tab:capops}
\begin{tabularx}{\linewidth}{@{}l X l@{}}
\toprule
\textbf{Instruction} & \textbf{Result} & \textbf{Source afterwards}\\\midrule
\insn{cmov cd, cs} & same capability in \reg{cd} & empty if linear\\
\insn{crestrict cd, cs, rm} & permissions ANDed with \texttt{rm<5:0>} & empty if linear\\
\insn{csplit cd, cs, rx} & \reg{cs} keeps $[b, b+x)$, \reg{cd} gets $[b+x, b+l)$ & shortened\\
\insn{cjoin cd, cs} & \reg{cd} extended by \reg{cs}; \reg{cs} must start where \reg{cd} ends & empty\\
\insn{cshare cd, cs} & perms $(\perm{R}\cup\perm{X})\cap p$, plus \perm{D} & empty if linear\\
\insn{catomic cd, cs} & perms exactly \perm{A}\perm{D}; source must be writable T2/T3 & empty\\
\insn{cclear cd} & empty & ---\\
\insn{cdom cd, cs} & \reg{cd}: revocable working capability for the region & revoker (\perm{V})\\
\insn{crevoke cd, cs} & \reg{cd}: fresh working capability; all older ones dead & revoker kept\\
\insn{cundom cd, cs} & \reg{cd}: plain capability for the region; all older ones dead & empty\\
\insn{cmkep cd, rx} & \reg{cd}: endpoint for node \reg{rx} & ---\\
\bottomrule
\end{tabularx}
\end{table}

Every derived capability keeps the revocable flag and generation of its source, so everything cut
from a working capability dies when the granules it came from are revoked.

\subsection{Sharing}
\insn{cshare} turns a capability into a non-writable, copyable one, keeping \perm{R} and \perm{X}. The
usual use is to publish data that will never change again: write it through a linear \perm{RW}
capability, then \insn{cshare} it and hand copies to every core. The boot ROM does exactly this with the
OS image (Section~\ref{sec:boot}). \insn{cshare} first waits until none of the core's memory
transfers is in flight, so every write through the writable capability has landed before any copy exists.

For memory that is not revocable, sharing is \emph{irrevocable}: the writable capability is
consumed and can never be reconstructed, because no operation adds \perm{W}. For revocable memory,
the holder of the revoker ends the sharing with \insn{crevoke} or \insn{cundom}, which kills every
share and returns a writable capability.

\subsection{Atomic segments}
\insn{catomic} turns a writable memory capability into an atomic-only, copyable one. Any number of
cores may hold it, and they can access the segment only with \insn{tamo}, which is indivisible. This
is how shared mutable state, such as a counter, a lock word or a mailbox, is created. The conversion
first waits for the core's in-flight memory transfers. Like sharing, it is irrevocable for memory
that is not revocable and ends with a revocation otherwise.

\subsection{Splitting and joining}
\insn{csplit} at offset $x$ requires $0 \le x \le \textit{len}$ and $x$ a multiple of 8; either piece may
be empty. \insn{cjoin}
is its inverse and checks that the two pieces are adjacent \emph{in order}, in the same space, with
identical permissions, the same revocable flag and the same generation, and, for revocable pieces,
that \reg{cs} does not begin a unit (Section~\ref{sec:domains}): its base is not the first byte of a
granule whose head bit is set. So pieces of two revocable regions never join, even when their
generations happen to be equal. A revoker (\perm{V}) is cut only at multiples of the granule size
(256~KiB); \insn{csplit} traps with \cause{capop} otherwise. Together with
\insn{cgive} and \insn{ctake}, these express data-parallel work division: split an array into one piece
per core, give the pieces away, take them back in core order and join them (the \texttt{parsum} and
\texttt{mandel} examples).

\subsection{Alignment}\label{sec:capalign}

\begin{rule*}{Capabilities are 8-byte aligned}
The base of every capability is a multiple of 8. \insn{csplit} traps with \cause{capop} (\csr{tval} =
the offset) unless the split offset is a multiple of 8, and no other instruction makes a base.
\end{rule*}

The reset capabilities start at 0; \insn{cjoin}, \insn{crestrict}, \insn{cshare}, \insn{catomic},
\insn{cdom}, \insn{crevoke} and \insn{cundom} keep a base; \insn{cmkep} makes an endpoint, whose base is a
node number. The rule is the same in every space. In memory it guarantees that an aligned offset
reaches exactly the doubleword at \texttt{base + offset}: a transfer at offset $o$ of a capability
touches the bytes from $\textit{base} + o$ and no others. In device space, ports, registers and
windows are 8-byte aligned (Chapter~\ref{ch:platform}), so a capability can be cut to exactly one
register, and transfers still reach single device bytes by offset. Lengths are not constrained: a
capability for the one-byte console output port has length 1.

\section{Revocation by address}\label{sec:domains}

Revocation lets the holder of a region of memory lend it out, in any number of pieces, shares and
atomic segments, to any number of cores, and take all of it back with one instruction. It works by
address: T3 and T2 are divided into \term{granules} of 256~KiB (T3, 64~MiB, has 256; T2, 256~KiB, has
one), and the machine keeps, for every granule, an entry shared by all cores (Table~\ref{tab:domtab}).

\begin{table}[htbp]
\centering\small
\caption{A granule's entry}\label{tab:domtab}
\begin{tabularx}{\linewidth}{@{}l l X@{}}
\toprule
\textbf{Field} & \textbf{Size} & \textbf{Meaning}\\\midrule
\texttt{GEN[g]} & 16 bits & the current generation; a revocable capability based in the granule and made under another one is dead\\
\texttt{HEAD[g]} & 1 bit & a \term{unit} begins here (below)\\
\bottomrule
\end{tabularx}
\end{table}

A capability is \term{revocable} or not. A revocable capability carries the generation its base's
granule had when it was made, and it is live only while the granule still has it
(Section~\ref{sec:live}). Revoking a range of granules advances their generations, which kills every
revocable capability \emph{based} in the range, however it was derived and wherever it is. Capabilities
derived from a revocable one never reach outside the range it was made for: they are cut from it, and
\insn{cjoin} never joins across units.

The granules a \insn{cdom} or a revocation covers as one form a \term{unit}: the first granule of
every unit has its head bit set, and so has the granule just after the end of a range \insn{cdom}
made revocable. Units do not change what is live; they bound what is revoked together and what may be
joined. Whenever a range is revoked, every unit it touches is revoked whole, because a capability of
that unit may reach into the range from outside it.

\begin{description}[leftmargin=1em,style=nextline]
\item[\insref{cdom} \texttt{cd, cs}] takes a live, writable, linear, non-revocable T2 or T3 capability
  whose base and length are multiples of 256~KiB (and the length not 0), and makes its range one
  unit: every granule of it gets the generation one past the highest any of them had, so no older
  capability can match it. \reg{cd} receives the \term{working capability} (the same segment and
  permissions, revocable, with that generation); \reg{cs} becomes the \term{revoker}, a linear,
  non-revocable capability for the same segment with permission \perm{V} only, which remembers the
  working capability's permissions. The revoker grants no access to memory.
\item[\insref{crevoke} \texttt{cd, cs}] revokes the revoker's range: every unit it touches gets the
  generation one past the highest among them, the range becomes a unit of its own, and so does what
  is left of each touched unit on either side. Every revocable capability based in those units, in any
  register, save table, inbox or device, is dead from this instruction on. \reg{cd} receives a fresh
  working capability for the revoker's range, with the permissions the revoker remembers; the revoker
  remains.
\item[\insref{cundom} \texttt{cd, cs}] revokes in the same way, consumes the revoker and puts a plain
  (non-revocable) capability for the range in \reg{cd}. The memory is back to ordinary use.
\end{description}

The three are privileged. A supervisor that runs a program in a revocable region lends it the working
capability and keeps the revoker in its save table; whatever the program does with the memory (split
it, share it, make parts atomic, give pieces to other cores), \insn{crevoke} or \insn{cundom} recovers
all of it when the program ends.

\begin{rule*}{Revokers}
A revoker is a linear capability with \perm{V} and is never revocable itself: no revocation kills it,
and \insn{cundom} consumes it. It is cut (\insn{csplit}) only at multiples of 256~KiB and joined
(\insn{cjoin}) only with an adjacent revoker that remembers the same permissions. A piece revokes its
own range: a supervisor can keep one revoker for a large region and revoke any granule-aligned part
of it by cutting the revoker there, revoking the piece and joining the pieces again. Because the
units a revocation touches are revoked whole, the rule that at most one live writable capability
covers any byte (Section~\ref{sec:linear}) holds.
\end{rule*}

The number of revocable regions is limited only by memory. A generation is 16 bits wide: a dead
capability could come back to life once its granule's generation has advanced through all 65,536
values while the capability is still held. A revocation of one unit advances its generation by one; a
revocation or \insn{cdom} that spans granules with different generations advances them all to one
past the highest, which can advance a granule by more than one.

A revocable capability's liveness is decided by the granule holding its base. A piece of length zero
cut at the very end of a range therefore has its base in the next granule: it lives and dies with
that granule (it grants no access in either case), and a revocable capability based past the end of
its space is never live.

\section{Endpoints}\label{sec:endpoints}

An \term{endpoint} is a capability in space 6 that names a node: its \textit{base} is the physical node
number. It grants no access to memory, and every transfer through it traps with \cause{perm}. It is used
with two instructions:

\begin{description}[leftmargin=1em,style=nextline]
\item[\insref{csend} \texttt{cs, ce}] gives \reg{cs} to the node the endpoint names, exactly like
  \insn{cgive} (same queues, same publication rule), but the node is not relative to the core window
  and not limited by it. Requires \perm{W} on the endpoint.
\item[\insref{crecv} \texttt{cd, ce}] takes the oldest gift from the node the endpoint names, like
  \insn{ctake cd, rs}. Requires \perm{R} on the endpoint; without it the instruction traps with
  \cause{capinv} (\csr{tval} = the endpoint's register number).
\end{description}

Both execute in user mode. \insn{ctake}, \insn{cpoll} and the inbox bit of \insn{tany} do not
(Section~\ref{sec:userinbox}): a user program receives only through an endpoint with \perm{R}, from
exactly the node that endpoint names, and only if its supervisor chose to give it one.

Only the supervisor can make an endpoint (\insref{cmkep}), with permissions \perm{R}\perm{W}\perm{D}.
\insn{crestrict} can reduce it to send-only or receive-only. Endpoints are how a supervisor lets a user
program, confined to its core window, talk to exactly the nodes it chooses. StrataOS gives every
process send-only endpoints to the cores of its servers and no receive rights at all: a request is a
message buffer sent through such an endpoint, and the reply is delivered by the kernel of the
process's core (Section~\ref{sec:os-msg}).

\section{Where capabilities live}

\begin{center}\small
\begin{tabularx}{\linewidth}{@{}l l X@{}}
\toprule
\textbf{Location} & \textbf{Entries} & \textbf{Accessed by}\\\midrule
capability registers & 16 per core & all capability instructions\\
save table & $\textit{CTAB}$ per core (256 by default, at least 64: the platform table says) & \insn{csave}, \insn{crestore} (slots 0--63), \insn{csavex}, \insn{crestorex} (every slot; supervisor only)\\
inbox & 2 per sender per core; 4 at the block device & \insn{cgive}, \insn{csend} (append); \insn{ctake}, \insn{crecv} (remove); \insn{cpoll}, \insn{tany} (inspect)\\
block device & while it performs a command & returned to a core's inbox on completion\\
\bottomrule
\end{tabularx}
\end{center}

A capability disappears for good when it is cleared, when the core holding it halts, or when it is
overwritten as a copyable or dead capability. A revocable capability becomes dead, wherever it is, when
the granule its base lies in is revoked. Memory whose only linear capability disappears is unreachable
for the rest of the run, unless it is revocable and its revoker survives.
